\documentclass[border=0pt]{standalone}
% The shared preamble of every docs diagram (docs/diagrams/<name>.tex).
% dvisvgm's driver first, so patterns and opacity survive into the SVG.
\def\pgfsysdriver{pgfsys-dvisvgm.def}
\usepackage[siunitx,RPvoltages]{circuitikz}
\usepackage{xcolor}
% Latin Modern: outline fonts for every glyph a label uses, µ included.
\usepackage[T1]{fontenc}
\usepackage{lmodern}
\usetikzlibrary{arrows.meta,backgrounds,calc,patterns,positioning}
\tikzset{
  % A white card, so diagrams read in mdBook's dark themes too.
  every picture/.append style={show background rectangle,
    background rectangle/.style={fill=white,inner sep=4pt}},
  % Layout geometry: layers told apart by fill and hatching, never colour alone.
  layer a/.style={draw=black,line width=0.6pt,fill=gray!20},
  layer b/.style={draw=black,line width=0.6pt,fill=gray!45},
  layer c/.style={draw=black,line width=0.6pt,fill=white,postaction={pattern=north east lines}},
  layer d/.style={draw=black,line width=0.6pt,fill=white,postaction={pattern=crosshatch dots}},
  marker/.style={draw=black,dashed,line width=0.6pt},
  % A measured distance and its label.
  dim/.style={{Latex[length=2mm]}-{Latex[length=2mm]},line width=0.5pt},
  dim label/.style={font=\small,fill=white,inner sep=1pt},
  % Where a rule fails: thick, dashed and labelled, not only red.
  violation/.style={draw=red!75!black,line width=1.4pt,dash pattern=on 3pt off 2pt},
  violation label/.style={font=\small\bfseries,text=red!75!black},
  caption/.style={font=\small},
}

\begin{document}
\begin{circuitikz}[font=\small]
  % Left: the reference netlist, an inverter whose n-channel side is two
  % cards in parallel (M1, M3). Right: the layout, one n-channel transistor.
  \node[font=\bfseries] at (1.5,5.0) {Reference netlist};
  \node[pmos] (p) at (1.5,2.8) {};
  \node[right=2pt] at (p.east) {M2 pfet, \SI{2}{\micro\metre}};
  \draw (p.S) -- ++(0,0.3) node[vcc] {VDD};
  \node[nmos] (n1) at (0.9,0.6) {};
  \node[nmos, xscale=-1] (n3) at (2.6,0.6) {};
  \node[left=1pt,font=\footnotesize] at (n1.south) {M1};
  \node[right=4pt,font=\footnotesize] at (n3.south) {M3};
  \draw (n1.D) -- ++(0,0.3) coordinate (d1) -- (d1 -| n3.D) coordinate (d3) -- (n3.D);
  \draw (p.D) |- (d1);
  \draw (d3) -- ++(2.0,0) node[right] {Y};
  \draw (n1.S) -- ++(0,-0.25) coordinate (s1) -- (s1 -| n3.S) coordinate (s3) -- (n3.S);
  \draw ($(s1)!0.5!(s3)$) node[ground] {};
  \draw (p.G) -- ++(-1.3,0) coordinate (a) -- (a |- n1.G) -- (n1.G);
  \draw (n3.G) -- ++(0.3,0) coordinate (g3) -- ++(0,-1.75) coordinate (g3b) -- (g3b -| a) -- (a |- n1.G);
  \draw (a |- p.G) node[left] {A};
  \draw[violation] (-1.0,-1.4) rectangle (4.25,1.95);
  \node[anchor=north west] at (-1.0,-1.45) {M1, M3: \SI{1}{\micro\metre} each};

  \node[font=\bfseries] at (7.5,5.0) {Layout, extracted};
  \node[pmos] (lp) at (7.5,2.8) {};
  \node[right=2pt] at (lp.east) {pfet, \SI{2}{\micro\metre}};
  \draw (lp.S) -- ++(0,0.3) node[vcc] {VDD};
  \node[nmos] (ln) at (7.5,0.6) {};
  \node[right=2pt] at (ln.east) {nfet, \SI{1}{\micro\metre}};
  \draw (lp.D) -- (ln.D);
  \draw (lp.D) ++(0,-0.35) coordinate (ly) -- ++(1.4,0) node[right] {Y};
  \draw (ln.S) node[ground] {};
  \draw (lp.G) -- ++(-0.6,0) coordinate (la) -- (la |- ln.G) -- (ln.G);
  \draw (la |- lp.G) node[left] {A};
  \draw[violation] (6.3,-0.9) rectangle (9.6,1.5);
  \node[violation label, align=center] at (4.8,-2.3) {w after merging: \SI{2}{\micro\metre} against \SI{1}{\micro\metre}};
\end{circuitikz}
\end{document}
